\documentclass{article}
\usepackage{amsmath, amssymb, amsthm}
\usepackage{hyperref}
\usepackage[margin=1in]{geometry}

\newtheorem{theorem}{Theorem}

\title{\vspace{-2em}Formalizing Origami Geometry in Lean 4}
\author{Manya Kansal}
\date{}

\begin{document}
\maketitle
\vspace{-1em}

\section*{Introduction}

The Huzita-Justin axioms (HJAs) formalise the seven fundamental origami
fold operations. Unlike just compass-and-straightedge operations, origami operations can solve cubic
equations via axiom O6 (the Beloch fold) \cite{alperin2000}, enabling constructions such as
angle trisection. The field of origami-constructible numbers is
characterised by Alperin's theorem: $z \in \mathbb{R}$ is
origami-constructible iff it lies in an extension of $\mathbb{Q}$ of
degree $2^a3^b$. Despite the elegance of these results, no
machine-verified proof exists.

\section*{Goals}

\begin{enumerate}
    \item Define terms for folds (\texttt{FoldLine}, \texttt{foldPoint},
    \texttt{foldLine}) and the HJAs O1--O7 as a synthetic axiomatic
    system rather than a coordinate system, similar to the SystemE/LeanGeo library \cite{avigad2009}.

    \item Prove basic fold properties: isometry, injectivity, and that
    the O2 fold line is the perpendicular bisector of the two given points, etc.

    \item Formalize \textbf{Abe's angle trisection} \cite{fuchs2011}:
    express the construction in terms of HJA axioms, then prove
    correctness by reducing to the identity
    $\cos(3\varphi) = 4\cos^3(\varphi) - 3\cos(\varphi)$
    using Mathlib's trigonometry library.

    \item If time permits, characterize which regular $n$-gons are
    origami-constructible (those with $\phi(n) = 2^c3^d$).
\end{enumerate}

\section*{The Plan}

Use a synthetic approach where \texttt{Point} and \texttt{Line} are
opaque types, following SystemE \cite{avigad2009}. Fold operations are
new primitives with axioms: involution, isometry, and the fact that fold-line
points are fixed. The HJAs are existence statements, e.g.\ O2 asserts
$\exists f,\ \texttt{foldPoint}\ f\ P_1 = P_2$. Correctness of Abe's
construction follows the arithmetic proof of Fuchs \cite{fuchs2011},
translated into Lean using Mathlib.

\begin{thebibliography}{9}
\setlength{\itemsep}{0pt}

\bibitem{avigad2009}
J.~Avigad, E.~Dean, J.~Mumma,
``A formal system for Euclid's Elements,''
\textit{Review of Symbolic Logic}, 2009.

\bibitem{fuchs2011}
C.~Fuchs,
``Angle trisection with origami and related topics,''
\textit{Elemente der Mathematik}, 2011.

\bibitem{alperin2000}
R.~C.~Alperin, R.~J.~Lang,
``One-, two-, and multi-fold origami axioms,''
\textit{Origami}, 2009.

\end{thebibliography}

\end{document}